% \begin{abstract}
% High-quality 3D assets are crucial for many downstream industries like virtual reality and entertainment. Image-conditioned diffusion/flow models excel in fidelity but suffer from viewpoint bias when the input view is limited or ambiguous, whereas text-conditioned models provide stable semantic priors yet lack fine-grained detail. We investigate whether vision and language can act as complementary signals for 3D generation and observe that even a simple late fusion of image- and text-conditioned predictions yields notable gains—evidence of strong cross-modal complementarity. This motivates us to propose a new task named Joint Image–Text Conditioned 3D Generation (JIT-3D), which requires the model to jointly consider the image and text condition in the 3D generation phase. To solve this task, we present \textbf{TIGON}, a minimalist cross-modal framework that couples an image-conditioned model with a text-conditioned model via (i) zero-initialized $1\times1$ convolutional bridges that enable feature sharing without disturbing pretrained behavior and (ii) step-wise prediction averaging during denoising. This design improves stability and generalization, especially under low-information views, and further supports mixing and editing conditions across modalities for flexible 3D asset creation. Extensive experiments on standard benchmarks and targeted low-information stress tests show TIGON has consistent performance gain over strong image-only and text-only baselines, validating complementary vision–language guidance as an effective principle for 3D generation.
% \end{abstract}

\begin{abstract}
% High-quality 3D assets are critical for VR/AR, industrial design, and entertainment, driving growing interest in generative models that can create 3D content from user-provided prompts. Most existing 3D generators, however, rely on a single conditioning modality: image-conditioned models deliver high visual fidelity by exploiting pixel-aligned cues but suffer from viewpoint bias when the input view is limited or ambiguous, whereas text-conditioned models benefit from broad semantic guidance yet lack low-level visual detail. This restricts how users can express their intent and raises a natural question: can the two modalities be combined to yield more flexible and faithful 3D generation? Our diagnostic study shows that even a simple late fusion of text- and image-conditioned predictions improves over single-modality models, evidencing strong cross-modal complementarity. Accordingly, we formalize the task of \textbf{Text–Image Conditioned 3D Generation}, which requires joint reasoning over a visual exemplar and a textual specification during generation. To address this task, we introduce \textbf{TIGON}, a minimalist dual-branch baseline that maintains separate image- and text-conditioned backbones with lightweight cross-modal fusion. Extensive experiments demonstrate that text–image conditioning yields consistent gains over single-modality methods, suggesting complementary vision–language guidance as a promising direction for future 3D generation research.  Project page: \url{https://jumpat.github.io/tigon-project-page}
High-quality 3D assets are essential for VR/AR, industrial design, and entertainment, motivating growing interest in generative models that create 3D content from user prompts. Most existing 3D generators, however, rely on a single conditioning modality: image-conditioned models achieve high visual fidelity by exploiting pixel-aligned cues bhaodeut suffer from viewpoint bias when the input view is limited or ambiguous, while text-conditioned models provide broad semantic guidance yet lack low-level visual detail. This limits how users can express intent and raises a natural question: can these two modalities be combined for more flexible and faithful 3D generation? Our diagnostic study shows that even simple late fusion of text- and image-conditioned predictions outperforms single-modality models, revealing strong cross-modal complementarity. We therefore formalize \textbf{Text–Image Conditioned 3D Generation}, which requires joint reasoning over a visual exemplar and a textual specification. To address this task, we introduce \textbf{TIGON}, a minimalist dual-branch baseline with separate image- and text-conditioned backbones and lightweight cross-modal fusion. Extensive experiments show that text–image conditioning consistently improves over single-modality methods, highlighting complementary vision–language guidance as a promising direction for future 3D generation research. Project page: \url{https://jumpat.github.io/tigon-page}
\end{abstract}



% High-quality 3D assets are critical for VR/AR, industrial design, and entertainment. Most current 3D generators rely on a \emph{single} conditioning modality: image-conditioned models exploit pixel-aligned cues but struggle with unseen views, whereas text-conditioned models provide broad semantic guidance but lack low-level detail. This motivates a natural question: can the two modalities be combined to yield more flexible and faithful 3D synthesis?

% We first conduct a diagnostic study showing that even a simple late fusion of text- and image-conditioned predictions consistently outperforms single-modality baselines, evidencing strong cross-modal complementarity. Building on this observation, we formalize the task of \textbf{Text--Image Conditioned 3D Generation}, which requires joint reasoning over a visual exemplar and a textual specification during synthesis, and establish a unified evaluation protocol for this setting. 

% To instantiate this task, we introduce \textbf{TIGON}, a minimalist dual-branch baseline that maintains separate image- and text-conditioned backbones with lightweight cross-modal fusion, while supporting free-form conditioning (image-only, text-only, or joint). Extensive experiments on standard benchmarks demonstrate that text--image conditioning yields consistent gains over strong single-modality baselines, suggesting complementary vision--language guidance as a simple and effective principle for future 3D generation research.